% !TeX root = 01_intro_ml_polytech.tex
% !TeX program = pdflatex
% !TeX encoding = UTF-8
\documentclass[t,10pt,aspectratio=169,rounded]{beamer}
\usetheme{CEA23}

\usepackage[utf8x]{inputenc}
\usepackage[english]{babel}
\usepackage{csquotes}

% Code listings: Carbon-style "code snapshot" cards (see code_listings_demo.tex)
\usepackage{beamer-code-listings}

\AtBeginSection[]{\sepframe[\insertsectionnumber]{}{}}

\begin{document}

\title{Introduction to Machine Learning}
\subtitle{Module 1 — Foundations: Python, Statistics, and Exploratory Data Analysis}
\date[30/11/2026]{30 November 2026}
\author{Riccardo \textsc{Finotello}}

\titlepagelogo{\includegraphics[height=24mm]{graphics/ORIGINAL/ISAS.png}}
% TODO: activate once img/titlepage.png exists in slides/img/
% \titlepagepicture{img/titlepage.png}
\setvalue{\diffusion}{}

\seminar[PolytechParis]{Machine Learning - PolytechParis}

\CEAmerci{Thank you}
\CEAcenter{CEA SACLAY}
\addr{91\,191 Gif-sur-Yvette Cedex\\France}
\CEAemail{\href{mailto:riccardo.finotello@cea.fr}{\texttt{riccardo.finotello@cea.fr}}}

% --- TITLE SLIDE ---
\begin{frame}[plain]
	\titlepage
\end{frame}

% ------------------------------------------------------------------
% Content slides: drop your module slide fragments into 01_slides/
% and include them below, e.g.
%   \input{01_slides/section_01}
%   \input{01_slides/section_02}
% Each fragment contains \section{...} and \begin{frame}...\end{frame}
% only (no preamble, no \begin{document}).
% ------------------------------------------------------------------

% --- CONCLUDING SLIDE ---
\begin{frame}[plain]
	\usebeamertemplate*{last page}
\end{frame}
\end{document}
